% Gedeeld door de syllabi van labo 3, 4 en 5. Pas dit hier aan, niet in de labobestanden.
\section{Overzicht}

Tijdens labo 3, 4 en 5 ben jij de nieuwe data-analist van een
(fictieve) jeugdclub voor voetbal en atletiek. De hoofdtrainer bezorgt je
drie exportbestanden uit de systemen van de club. Die zijn bewust
rommelig, zoals echte exports: inconsistente labels, sprinttijden in
milliseconden, ontbrekende waarden, een atleet die twee keer voorkomt en
een kolom met een komma als decimaalteken.

Elke sessie bouwt verder op de vorige:

\begin{description}[leftmargin=1.6cm, style=nextline]
  \item[Labo 3] \textbf{Van bestand tot nette tabel.} Inladen, verkennen,
    filteren, kolommen toevoegen en opkuisen tot je de data kunt
    vertrouwen.
  \item[Labo 4] \textbf{Combineren, samenvatten, herschikken.} Groepen
    samenvatten, tabellen koppelen, wisselen tussen een brede en een lange
    vorm, en werken met een hartslagreeks per seconde.
  \item[Labo 5] \textbf{Vertel het verhaal met plotly.} De juiste grafiek
    kiezen voor een vraag, eerlijke grafieken maken met
    \code{plotly.express} en ze combineren tot een klein interactief
    dashboard.
\end{description}
